% Package definition
\usepackage[utf8]{inputenc}
\usepackage[T1]{fontenc}
\usepackage[sfdefault]{atkinson}  % Atkinson Hyperlegible font (mandatory UNamur)
\usepackage[english]{babel}
\usepackage[top=2cm, bottom=2.5cm, left=2cm, right=2cm]{geometry}
\usepackage{setspace}
\usepackage{csquotes}
\usepackage{acronym}

\renewcommand{\familydefault}{\sfdefault}
\usepackage{microtype}

\usepackage{graphicx}
\usepackage{float}
\usepackage[dvipsnames]{xcolor}
\usepackage{tikz}
\usepackage{fancybox}

\graphicspath{{figures/}{img/}}

\usepackage{amsmath}
\usepackage{amssymb}
\usepackage{amsthm}

\usepackage{array}
\usepackage{booktabs}
\usepackage{longtable}
\usepackage{multirow}
\usepackage{tabularx}
\usepackage{etoolbox}
\AtBeginEnvironment{tabular}{\singlespacing}
\AtBeginEnvironment{tabularx}{\singlespacing}

\usepackage{enumitem}
\usepackage{pifont}

\usepackage[backend=biber,style=authoryear,maxcitenames=2,maxbibnames=99,uniquename=false,uniquelist=false,dashed=false]{biblatex}
\addbibresource{bibliography.bib}

% \cite{} produces [Author, Year] with square brackets and comma
\renewcommand*{\mkbibparens}[1]{\mkbibbrackets{#1}}
\AtBeginDocument{\let\cite\parencite}
% Force comma between author and year
\renewcommand*{\nameyeardelim}{\addcomma\addspace}

% et al without trailing period
\DefineBibliographyStrings{english}{andothers = {et~al}}

% For entries with a shorthand (RFCs), cite as [RFC1034, 1987]
\renewbibmacro*{cite:shorthand}{%
  \printfield{shorthand}%
  \setunit{\nameyeardelim}%
  \printfield{labelyear}}

\usepackage[hidelinks]{hyperref}
\usepackage{url}
\usepackage{cleveref}

\hypersetup{
    colorlinks=false,
    linkcolor=black,
    citecolor=black,
    urlcolor=blue,
    pdftitle={DNS Measurements in Space and Time},
    pdfauthor={Olivier Gautier},
    pdfsubject={Master 60 Computer Science - UNamur},
    pdfkeywords={DNS, RIPE Atlas, distributed measurements, geolocation}
}

\urlstyle{same}

\usepackage{listings}
\usepackage[linesnumbered,ruled,vlined]{algorithm2e}

% Configuration listings - Couleurs
\definecolor{codebg}{rgb}{0.97,0.97,0.97}
\definecolor{javakeyword}{RGB}{0,0,180}
\definecolor{javacomment}{RGB}{0,128,0}
\definecolor{javastring}{RGB}{163,21,21}

% Style Python
\lstdefinestyle{python}{
  language=Python,
  backgroundcolor=\color{codebg},
  basicstyle=\ttfamily\small,
  keywordstyle=\color{blue}\bfseries,
  commentstyle=\color{gray}\itshape,
  stringstyle=\color{orange},
  showstringspaces=false,
  numbers=left,
  numberstyle=\tiny\color{gray},
  numbersep=10pt,
  frame=single,
  tabsize=4,
  captionpos=b,
  breaklines=true,
  breakatwhitespace=true
}

% Style Bash
\lstdefinestyle{bash}{
  language=bash,
  backgroundcolor=\color{codebg},
  basicstyle=\ttfamily\small,
  keywordstyle=\color{blue}\bfseries,
  commentstyle=\color{gray}\itshape,
  stringstyle=\color{orange},
  showstringspaces=false,
  numbers=left,
  numberstyle=\tiny\color{gray},
  numbersep=10pt,
  frame=single,
  breaklines=true
}

% Style JSON
\lstdefinelanguage{json}{
  morestring=[b]",
  morestring=[d]',
  literate=
   *{0}{{{\color{blue}0}}}{1}
    {1}{{{\color{blue}1}}}{1}
    {2}{{{\color{blue}2}}}{1}
    {3}{{{\color{blue}3}}}{1}
    {4}{{{\color{blue}4}}}{1}
    {5}{{{\color{blue}5}}}{1}
    {6}{{{\color{blue}6}}}{1}
    {7}{{{\color{blue}7}}}{1}
    {8}{{{\color{blue}8}}}{1}
    {9}{{{\color{blue}9}}}{1}
    {:}{{{\color{red}:}}}{1}
    {,}{{{\color{red},}}}{1}
    {\{}{{{\color{red}\{}}}{1}
    {\}}{{{\color{red}\}}}}{1}
    {[}{{{\color{red}[}}}{1}
    {]}{{{\color{red}]}}}{1}
    {true}{{{\color{purple}true}}}{4}
    {false}{{{\color{purple}false}}}{5}
    {null}{{{\color{purple}null}}}{4},
}

\lstdefinestyle{json}{
  language=json,
  backgroundcolor=\color{codebg},
  basicstyle=\ttfamily\small,
  keywordstyle=\color{blue},
  stringstyle=\color{orange},
  showstringspaces=false,
  numbers=left,
  numberstyle=\tiny\color{gray},
  numbersep=10pt,
  frame=single,
  breaklines=true,
  captionpos=b
}

\usepackage{fancyhdr}

% \pagestyle{fancy} activé dans main.tex avant les chapitres
\fancyhf{}
\fancyhead[L]{\nouppercase{\leftmark}}
\fancyhead[R]{\thepage}
\renewcommand{\headrulewidth}{0pt}
\renewcommand{\footrulewidth}{0.4pt}
\renewcommand{\plainfootrulewidth}{0.4pt}

% Style pour pages de chapitre (première page)
\fancypagestyle{plain}{
    \fancyhf{}
    \fancyfoot[L]{\nouppercase{\leftmark}}
    \fancyfoot[R]{\thepage}
    \renewcommand{\headrulewidth}{0pt}
}

\setlength{\headheight}{13.59999pt}

% ============================================================================
% PACKAGES CAPTIONS
% ============================================================================

\usepackage{caption}
\usepackage{subcaption}

\captionsetup{
    font=small,
    labelfont=bf,
    textfont=it
}

% ============================================================================
% CONFIGURATION DOCUMENT
% ============================================================================

% Profondeur table des matières
\setcounter{tocdepth}{2}
\setcounter{secnumdepth}{3}

% Espacement paragraphes
\setlength{\parskip}{0.5em}
\setlength{\parindent}{0pt}

% Justification
\usepackage{ragged2e}

% ============================================================================
% COMMANDES PERSONNALISÉES
% ============================================================================

% Commande pour le symbole de référence
\newcommand{\refmark}[1]{\textsuperscript{\cite{#1}}}

% Convenience reference commands
\newcommand{\figref}[1]{Figure~\ref{#1}}
\newcommand{\tabref}[1]{Table~\ref{#1}}
\newcommand{\chapref}[1]{Chapter~\ref{#1}}
\newcommand{\secref}[1]{Section~\ref{#1}}